\documentclass[11pt]{article}

\usepackage[a4paper,margin=27mm]{geometry}
\usepackage{amsmath,amssymb,amsthm,mathtools}
\usepackage{booktabs}
\usepackage{enumitem}
\usepackage[hidelinks]{hyperref}

\newtheorem{definition}{Definition}[section]
\newtheorem{proposition}[definition]{Proposition}
\newtheorem{criterion}[definition]{Criterion}
\theoremstyle{remark}
\newtheorem{remark}[definition]{Remark}

\newcommand{\cC}{\mathcal C}
\newcommand{\cS}{\mathcal S}
\newcommand{\cE}{\mathcal E}
\newcommand{\cN}{\mathcal N}
\newcommand{\wt}{\operatorname{wt}}

\title{Frame-Aware Fault Tolerance for Transversal Phase Gates\\
\large Exact syndromes, noisy records, and a finite decoder test}
\author{}
\date{\today}

\begin{document}
\maketitle

\begin{abstract}
CSS fault-tolerant circuits often guarantee more structure than a bound on
the total support of the residual error: the two CSS components may each be
correctable even when their union is not.  A transversal phase gate changes
an $XZ$-structured error into a $YZ$-structured error, so decoding the two
CSS syndrome channels independently is no longer justified.  This note
gives circuit-level definitions of strict and frame-aware CSS fault
tolerance, proves that exact \emph{joint} syndrome decoding after a valid
transversal logical phase gate is always possible, and separates that
code-capacity statement from the harder problem of noisy syndrome
extraction.  A finite checker implementing the resulting ambiguity tests is
provided and evaluated on the $[[7,1,3]]$ Steane code and the $[[17,1,5]]$
colour code.
\end{abstract}

\paragraph{Main conclusion.}
There is no code-capacity counterexample under the intended hypotheses.  If
the incoming $XZ$ componentwise error set is correctable and the physical
transversal $S$ or $S^\dagger$ preserves the code, then the outgoing $YZ$
set is correctable by unitary conjugation.  Measuring all $Z$- and $X$-type
checks and decoding the two records jointly is therefore sufficient in the
ideal-syndrome model.  What remains nontrivial is the \emph{gadget-level}
question: faults in the two extraction circuits can corrupt both the data
and the records.  That enlarged, circuit-derived fault set must pass the
same finite ambiguity test.

\section{CSS notation and two notions of fault tolerance}

Let $\cC$ be an $[[n,k,d]]$ CSS code with stabiliser group $\cS$ and
$t=\lfloor(d-1)/2\rfloor$.  Let the rows of binary matrices $H_X$ and $H_Z$
specify its $X$- and $Z$-type stabiliser generators, so
$H_XH_Z^{\mathsf T}=0$.  If $A\subseteq\{1,\ldots,n\}$ and
$R\in\{X,Y,Z\}$, define
\[
 R_A:=\prod_{i\in A}R_i.
\]
Equivalently, for the indicator vector $a$ of $A$, one may write $R(a)$.
Thus the symbols $X_A$ and $Z_B$ are Pauli operators, while $A$ and $B$ are
sets of qubit positions.  For $E=X_AZ_B$, the complete CSS syndrome is
\[
 s(E)=\bigl(H_Za^{\mathsf T},\,H_Xb^{\mathsf T}\bigr),
 \tag{1}
\]
where the first component is the record from the $Z$-type checks and the
second is the record from the $X$-type checks.  Phases of Pauli operators
are ignored unless a stabiliser sign is relevant.

\subsection{Strict fault tolerance}

The definitions below apply to arbitrary gadgets, including logical gates,
error correction, measurement, and preparation.  Preparation is simply the
case with no encoded input.

\begin{definition}[Strict fault tolerance]
Consider a gadget $G$ with an ideal encoded operation $G_0$.  The gadget is
\emph{strictly $t$-fault tolerant} if, for all $s,r\geq0$ with
$s+r\leq t$, an input residual Pauli of total weight at most $s$ together
with at most $r$ internal faults produces, relative to $G_0$, an output
residual Pauli of total weight at most $s+r$.  For several blocks, weight is
counted globally across the participating blocks.  For a postselected
gadget the condition is imposed conditional on acceptance; erroneous
classical output bits are included in the same fault budget.
\end{definition}

This is the familiar ``errors plus faults do not proliferate'' condition
used in fault-tolerant gadget analyses \cite{AGP}.  It is deliberately a
property of the whole circuit, not only of its state-preparation subroutine.
For a preparation gadget, $s=0$ and the definition reduces to the statement
that $r$ accepted faults leave an error of weight at most $r$.

\subsection{Frame-aware CSS fault tolerance}

Let $F=(P,Q)$ be an ordered pair of distinct Pauli axes.  Define
\[
 \cE_F(r):=\{P_AQ_B: |A|\leq r,\ |B|\leq r\}.
 \tag{2}
\]
We call $\cE_F(r)$ the set of \emph{$F$-admissible errors of order $r$}.
This replaces the less precise language of an ``error promise''.  The
$F$-component weight of $P_AQ_B$ is
\[
 \wt_F(P_AQ_B):=\max(|A|,|B|).
\]
For a fixed ordered pair $F$, the sets $A$ and $B$ are unique up to the
irrelevant overall Pauli phase.

For several blocks with frames $F_j=(P_j,Q_j)$, use the global component
weight
\[
 \wt_{\boldsymbol F}\!\left(\prod_j(P_j)_{A_j}(Q_j)_{B_j}\right)
 :=\max\!\left(\sum_j|A_j|,\sum_j|B_j|\right).
 \tag{3}
\]
The global definition is important: allowing every input block to saturate
an independent weight-$t$ budget would permit a transversal two-block gate
to combine them into a weight-$2t$ component on one output.

\begin{definition}[Frame-aware CSS fault tolerance]
A gadget with prescribed input and output frames is
\emph{$\boldsymbol F$--CSS-$t$-fault tolerant} if, whenever the input has
global component weight at most $s$ and the gadget contains at most $r$
faults with $s+r\leq t$, its residual output error has component weight at
most $s+r$ in the prescribed output frames, up to a known Pauli-frame
update.  Postselection and classical outputs are treated as in the strict
definition.
\end{definition}

Strict fault tolerance bounds the union of the two component supports.
Frame-aware CSS fault tolerance bounds them separately.  For example,
$X_AZ_B\in\cE_{XZ}(r)$ can have total support $2r$.

\begin{proposition}[Why the $XZ$ model is correctable]
Let $d_X$ and $d_Z$ be the minimum weights of nontrivial pure-$X$ and
pure-$Z$ logical operators.  If $2t<\min(d_X,d_Z)$, then
$\cE_{XZ}(t)$ is a correctable error set.
\end{proposition}

\begin{proof}
For $E=X_AZ_B$ and $E'=X_{A'}Z_{B'}$ in $\cE_{XZ}(t)$, the pure-$X$ and
pure-$Z$ components of $E^\dagger E'$ each have weight at most $2t$.  If
$E^\dagger E'$ has zero syndrome, the CSS form implies that these two
components separately normalise $\cS$.  The pure-distance bounds force each
component to be a stabiliser.  Hence $E^\dagger E'$ is not a nontrivial
logical operator, and the Knill--Laflamme condition proves correctability
\cite{KnillLaflamme}.
\end{proof}

\section{What transversal \texorpdfstring{$S$}{S} changes}

Ignoring phases,
\[
 H:(X,Z)\mapsto(Z,X),\qquad
 S^{\pm1}:(X,Z)\mapsto(Y,Z),\qquad
 HSH:(X,Z)\mapsto(X,Y).
 \tag{4}
\]
Thus $H$ preserves the unordered $XZ$ frame, $S$ or $S^\dagger$ maps it to
$YZ$, and $HSH$ maps it to $XY$.

There are two sound responses.

\paragraph{Route A: use strict fault tolerance.}
A tensor product of single-qubit Clifford gates changes Pauli axes but not
their support.  Therefore a strictly fault-tolerant input remains within
the same total-weight bound after transversal $S$.  This is the original
strict-preparation solution, now viewed as the preparation instance of a
general circuit property.

\paragraph{Route B: retain the frame and decode jointly.}
The usual independent CSS decoder is justified for $XZ$ errors by the
previous proposition.  After $S$, the $Y$ component contributes to both
binary syndrome channels, so those two records must be interpreted as one
joint record.  This is not an appeal to a small maximum-likelihood
improvement; it is a change in the set that the decoder is required to
correct.

The central point is the following elementary theorem.

\begin{proposition}[The ideal joint decoder always exists]
Let $U$ be the physical transversal $S$ or $S^\dagger$ used by the protocol,
and assume that $U\cC=\cC$ (so $U$ implements a valid logical Clifford on
the code).  If $\cE_{XZ}(t)$ is correctable, then
\[
 \cE_U(t):=U\cE_{XZ}(t)U^\dagger
 \tag{5}
\]
is correctable on the same code.  For uniform transversal $S^{\pm1}$,
$\cE_U(t)=\cE_{YZ}(t)$ up to phases.
\end{proposition}

\begin{proof}
Conjugate the Knill--Laflamme equations, or a recovery channel for
$\cE_{XZ}(t)$, by $U$.  Because $U$ preserves the code space, the conjugated
recovery acts on the same code.
\end{proof}

Many two-dimensional colour codes implement the Clifford group
transversally \cite{BombinColor}; depending on conventions and code length,
the physical $S$ or $S^\dagger$ may realise logical $S$ or $S^\dagger$.
Only the code-preserving physical operation actually used by the protocol
matters here.  We do not consider codes for which the proposed transversal
phase operation has no valid logical action.

\begin{remark}[No ideal counterexample under the hypotheses]
An ideal-syndrome counterexample consisting of two post-$S$ errors with the
same syndrome and different logical actions cannot exist under the
proposition's assumptions.  Such a pair would violate correctability of
$U\cE_{XZ}(t)U^\dagger$, and conjugating it back would violate
correctability of $\cE_{XZ}(t)$.  A counterexample must therefore leave the
hypotheses: the physical $S$ does not preserve the code, the incoming fault
set is larger than the certified $XZ$ set, or the syndrome record is noisy
or incomplete.
\end{remark}

\section{Joint syndrome measurement and frame reset}

\begin{proposition}[Two CSS channels are the full syndrome]
Suppose every $Z$-type and every $X$-type stabiliser generator is measured
exactly, and no data fault occurs between the two groups of measurements.
Measuring the $Z$-type checks and then the $X$-type checks, or in the reverse
order, returns precisely the full syndrome in Eq.~(1).  Hence a joint lookup
table on the two records implements an ideal decoder for $\cE_U(t)$.
\end{proposition}

\begin{proof}
All stabilisers commute.  The first group returns $H_Zx^{\mathsf T}$ and the
second returns $H_Xz^{\mathsf T}$ for the same data Pauli $X(x)Z(z)$.  Their
concatenation is $s(E)$, independent of measurement order.
\end{proof}

This is the clean, circuit-independent statement: ideal frame reset is
nothing more than obtaining both CSS syndrome channels and decoding them
together.  Steane error correction is one physical way to acquire these
records.  It couples verified encoded $|0_L\rangle$ and $|+_L\rangle$
ancillas transversally to the data \cite{SteaneEC}.  Other extraction
circuits may provide the same information.

The ideal equivalence does \emph{not} make all implementations equivalent
under faults.  A fault between the two groups of measurements changes the
data seen by the second group; an ancilla fault can both propagate to the
data and corrupt a record; and a classical readout fault changes the record
without changing the data.  These events must be analysed using the actual
extraction circuit.

\subsection{Noisy records: two different questions}

Let a candidate fault event be a pair $(E,m)$, where $E$ is its effective
data Pauli and $m$ is a binary error in the measured syndrome bits.  The
decoder observes
\[
 \widetilde s(E,m)=s(E)\oplus m.
 \tag{6}
\]
Let $\cN(\cS)$ denote the Pauli normaliser of the stabiliser group
$\cS$; the quotient $\cN(\cS)/\cS$ is the logical Pauli group.

\begin{criterion}[Logical ambiguity]
For every pair of candidate events with the same observed record, require
\[
 \widetilde s(E,m)=\widetilde s(E',m')
 \quad\Longrightarrow\quad
 E^\dagger E'\notin\cN(\cS)\setminus\cS.
 \tag{7}
\]
If this fails, the record is compatible with two events differing by a
nontrivial logical Pauli.  No correlated decoder or subsequent recovery can
correct both branches.
\end{criterion}

This is the Pauli form of the Knill--Laflamme condition applied to the
conditional error set associated with one observed record.  It captures the
fundamental question raised in the discussion.

\begin{criterion}[Exact one-shot reset]
A single Pauli correction chosen only from the observed record returns every
candidate exactly to the same codespace state only if
\[
 \widetilde s(E,m)=\widetilde s(E',m')
 \quad\Longrightarrow\quad E^\dagger E'\in\cS.
 \tag{8}
\]
\end{criterion}

With perfect syndrome bits, equal observed records imply equal true
syndromes, so Eqs.~(7) and (8) are equivalent.  With measurement errors they
differ.  Two events can share an observed record while $E^\dagger E'$ has a
nonzero syndrome.  This prevents an exact one-shot reset, but it is a
detectable residual error rather than an immediate logical ambiguity.  A
fault-tolerance proof must then show that all such residuals stay within the
allowed output error set.

Code perfection does not remove either requirement.  In particular, a
nondegenerate perfect code already assigns all available syndromes to
correctable data errors; admitting independent errors in the syndrome bits
creates competing data-versus-measurement explanations rather than
eliminating them.

\subsection{Destructive measurement and injection}

A destructive $Z$-basis measurement is naturally compatible with an $XZ$
or $YZ$ frame: only one frame component ($X$ or $Y$) flips its outcomes.  In
an $XY$ frame both bounded components can flip outcomes, so their union can
exceed the classical decoder's radius.  Measuring the complementary CSS
syndrome first and decoding it jointly with the destructive record supplies
the same two-channel information as a full ideal correction cycle.

The same accounting applies to the usual $T$-state injection circuit in
which the resource controls a CNOT into the input and the input is then
measured destructively in $Z$.  The CNOT accepts an input target frame
containing $X$ ($XZ$ or $XY$), while the measurement accepts a frame
containing $Z$ ($XZ$ or $YZ$).  Their intersection is $XZ$.  For a $YZ$ or
$XY$ input, the injection record supplies one syndrome channel and a
complementary extraction can supply the other.  Whether the resulting noisy
gadget is fault tolerant is again decided by Eqs.~(7) and (8), together with
the required residual-weight bound.

\section{A finite checker and concrete results}

The script
\begin{center}
\path{spiderstate/correlated_frame_decoder.py}
\end{center}
implements both criteria for any CSS code accepted by
\texttt{spiderstate.utils.load\_qecc}.  Its current code-capacity/readout
model is deliberately explicit:
\begin{enumerate}[leftmargin=*,itemsep=2pt]
  \item enumerate every pre-$S$ error $X_AZ_B$;
  \item assign it CSS component order $\max(|A|,|B|)$;
  \item propagate it through $S$ as
        $X_AZ_B\mapsto Y_AZ_B=X_AZ_{A\triangle B}$, up to phase;
  \item optionally enumerate a syndrome-bit error $m$, with total order
        $\max(|A|,|B|)+|m|\leq t$;
  \item group events by the observed joint record $\widetilde s$;
  \item count violations of the logical-ambiguity and exact-reset criteria.
\end{enumerate}
The script also checks, including stabiliser signs, that uniform transversal
$S$ preserves the loaded code.  Correlated data-and-record faults from a
specific extraction circuit are outside this simple model; they should be
fed into the same grouping test once enumerated.

The implementation was tested in the repository's \texttt{zxlive}
environment with
\begin{center}
\small\texttt{conda run -n zxlive python -m unittest
spiderstate.test\_correlated\_frame\_decoder -v}.
\end{center}
Both automated tests pass.  The exhaustive results are:

\begin{center}
\begin{tabular}{@{}llrrrr@{}}
\toprule
Code & record model & candidates & records & logical & exact reset \\
\midrule
$[[7,1,3]]$  & exact & 64    & 64    & 0 & 0 \\
$[[17,1,5]]$ & exact & 23716 & 13456 & 0 & 0 \\
$[[7,1,3]]$  & noisy & 70    & 64    & 0 & 6 \\
$[[17,1,5]]$ & noisy & 29020 & 14074 & 0 & 3018 \\
\bottomrule
\end{tabular}
\end{center}
Here ``logical'' and ``exact reset'' count ambiguous observed-record
buckets, not error pairs.  The exact-syndrome rows confirm the general
conjugation theorem.  The noisy rows show that the tested model creates no
logical ambiguity through order $t$, but a single Pauli selected from one
noisy record cannot always perform an exact reset.

For the $[[7,1,3]]$ code, one explicit order-one exact-reset collision is:
\begin{align*}
 &\text{event A: } E=I,\quad m=(1,0,\ldots,0),\\
 &\text{event B: } E_{\rm before}=X_1Z_1,\quad
   E_{\rm after\ }S=X_1,\quad m=0.
\end{align*}
The two events have the same observed record.  Their data errors differ by
$X_1$, which is detectable rather than logical.  This example starts with
an admissible $XZ$-structured first-order error and explicitly includes the
transversal $S$, as required.  It shows why noisy extraction is not an exact
code-capacity reset while also showing why that collision alone is not a
logical failure.

\section{Propagating frames through a Clifford circuit}

For completeness, the useful two-block rules follow directly from Pauli
propagation.  Under transversal CNOT, the control frame must contain $Z$ and
the target frame must contain $X$:
\[
 F_{\rm control}\in\{XZ,YZ\},\qquad
 F_{\rm target}\in\{XZ,XY\}.
 \tag{9}
\]
An $X$ or $Y$ component on the control propagates $X$ to the target; a $Z$
or $Y$ component on the target propagates $Z$ to the control.  Under
controlled-$Y$, the analogous sufficient conditions are
\[
 F_{\rm control}\in\{XZ,YZ\},\qquad
 F_{\rm target}\in\{YZ,XY\}.
 \tag{10}
\]
These are axis-compatibility rules.  The global component budget in
Eq.~(3) is still required to prevent two independently saturated inputs
from accumulating on one output block.

For a Clifford network whose gates obey these rules, the frame label records
how the certified input error set is propagated.  The decoder need not know
every earlier gate location in order to interpret an \emph{exact} full
syndrome, provided the propagated set is determined by the current frame.
For a noisy correction gadget, however, the decoder must additionally know
the effective data-and-record events generated by faults in that gadget.

\section{What is proved, and what remains}

\begin{enumerate}[leftmargin=*,itemsep=3pt]
  \item \textbf{Proved for arbitrary gadgets:} strict and frame-aware CSS
        fault tolerance are circuit properties; preparation is a special
        case.
  \item \textbf{Proved for exact syndromes:} after a code-preserving
        transversal logical $S$ or $S^\dagger$, the transformed admissible
        set is correctable.  Jointly measuring the two CSS syndrome channels
        is sufficient, and no logical collision can occur within that set.
  \item \textbf{Checked exhaustively:} the exact joint decoder passes for
        the $[[7,1,3]]$ Steane and $[[17,1,5]]$ colour codes through their
        full correction radii.
  \item \textbf{Also checked in a simple noisy-record model:} neither code
        has a logical ambiguity through order $t$, but both have one-shot
        exact-reset ambiguities.
  \item \textbf{Still gadget dependent:} a complete fault-tolerance proof
        for Steane EC, a partial extraction followed by destructive
        measurement, or magic-state injection must enumerate the actual
        propagated data errors and record corruptions.  The finite test in
        Eq.~(7) then decides logical ambiguity, while the desired strict or
        frame-aware output bound decides whether detectable residuals are
        acceptable.
\end{enumerate}

The conceptual payoff is that the ideal correlated-decoding question is
settled positively and generally.  The remaining work is concrete rather
than philosophical: derive the fault table of the chosen extraction gadget
and pass it through the same record-by-record logical-coset test.

\begin{thebibliography}{9}

\bibitem{KnillLaflamme}
E.~Knill and R.~Laflamme,
``A theory of quantum error-correcting codes,''
\emph{Physical Review A} \textbf{55}, 900 (1997),
\href{https://arxiv.org/abs/quant-ph/9604034}{arXiv:quant-ph/9604034}.

\bibitem{SteaneEC}
A.~M.~Steane,
``Active stabilisation, quantum computation and quantum state synthesis,''
\emph{Physical Review Letters} \textbf{78}, 2252 (1997),
\href{https://arxiv.org/abs/quant-ph/9611027}{arXiv:quant-ph/9611027}.

\bibitem{AGP}
P.~Aliferis, D.~Gottesman, and J.~Preskill,
``Quantum accuracy threshold for concatenated distance-3 codes,''
\emph{Quantum Information \& Computation} \textbf{6}, 97--165 (2006),
\href{https://arxiv.org/abs/quant-ph/0504218}{arXiv:quant-ph/0504218}.

\bibitem{BombinColor}
H.~Bombin and M.~A.~Martin-Delgado,
``Topological quantum distillation,''
\emph{Physical Review Letters} \textbf{97}, 180501 (2006),
\href{https://arxiv.org/abs/quant-ph/0605138}{arXiv:quant-ph/0605138}.

\end{thebibliography}

\end{document}
